% Book 2 Introduction
% Brief review and transition from Book 1

\chapter*{Introduction}
\addcontentsline{toc}{chapter}{Introduction}

This book is the second volume in a two-book series on Mathematical Programming and Operations Research.

\section*{Prerequisites}

This book assumes familiarity with the material covered in \textbf{Book 1: Introduction to Mathematical Programming}, including:

\begin{itemize}
    \item \textbf{Linear Programming (LP):} Formulating LP problems, the simplex method, duality theory, and sensitivity analysis
    \item \textbf{Discrete Algorithms:} Graph theory fundamentals, shortest path algorithms, and network flow concepts
    \item \textbf{Introduction to Integer Programming:} Basic IP formulations including knapsack problems, assignment problems, and set covering
\end{itemize}

\section*{What This Book Covers}

\textbf{Part I: Integer Programming} provides a comprehensive treatment of integer programming, building on the introductory material from Book 1. Topics include:
\begin{itemize}
    \item Advanced IP formulations and modeling techniques
    \item Branch-and-bound and cutting plane algorithms
    \item Exponential-sized formulations and column generation
    \item Computational complexity theory (P, NP, NP-complete)
    \item Heuristics and metaheuristics for combinatorial optimization
\end{itemize}

\textbf{Part II: Nonlinear Programming} introduces optimization with nonlinear objective functions and constraints:
\begin{itemize}
    \item Unconstrained optimization: optimality conditions, gradient descent, Newton's method
    \item One-dimensional optimization methods
    \item Gradient methods: steepest descent, conjugate gradient
    \item Constrained optimization: KKT conditions, barrier methods
    \item Applications to machine learning and data science
\end{itemize}

\section*{How to Use This Book}

Each chapter builds upon previous material. Readers should have a solid understanding of linear algebra (matrices, vectors, eigenvalues) and multivariable calculus (gradients, Hessians) for the nonlinear programming sections.

The appendices provide a review of essential linear algebra concepts for reference.

\vspace{1cm}
\begin{center}
\rule{0.5\textwidth}{0.4pt}
\end{center}
